%=====================================================================
\chapter{Measurement Theory and Instrument Validity}\label{ch:measurement}
%=====================================================================
\locator{4}{Part~4 --- Measurement, Data and Analysis}

\begin{outcomes}
By the end of this chapter you will be able to:
\begin{tight}
  \item \textbf{Classify} a measure by scale type and say what that permits.
  \item \textbf{Distinguish} the forms of validity and name a threat to each.
  \item \textbf{Compute} and \textbf{interpret} a reliability coefficient, and
    explain why Cronbach's alpha is weaker evidence than it appears.
  \item \textbf{Derive} measures from goals using GQM rather than measuring
    what is convenient.
  \item \textbf{Defend} a proxy measure, or abandon it.
\end{tight}
\end{outcomes}

\begin{prereq}
\chref{questions} for operationalisation. This chapter is where that becomes a
technical matter with tests attached. \chref{survey} if you are building an
instrument.
\end{prereq}

\begin{keyterms}
construct $\bullet$ measure $\bullet$ nominal, ordinal, interval, ratio
$\bullet$ face validity $\bullet$ content validity $\bullet$ construct validity
$\bullet$ convergent and discriminant validity $\bullet$ criterion validity
$\bullet$ reliability $\bullet$ test--retest $\bullet$ internal consistency
$\bullet$ Cronbach's alpha $\bullet$ McDonald's omega $\bullet$ GQM $\bullet$
proxy measure
\end{keyterms}

\section{The Weakest Link}

A study is only as good as its measures. This is stated often and acted on
rarely, because measurement is unglamorous and its failures are invisible in
the output: a regression on a bad measure produces coefficients, standard
errors and p-values exactly as it would on a good one. Nothing complains.

\begin{alertbox}[Analysis cannot repair measurement]
No statistical technique recovers a construct that was never captured. If your
measure of ``code quality'' is comment density, then a sophisticated
mixed-effects model on comment density tells you about comment density. The
model is not wrong; the inference from it to code quality is.

\medskip
This is why \reg{PhD Regs 2024}{\S14.3(B)} places validity --- ``does it measure
what it says?'' --- alongside reliability as a criterion your dissertation is
assessed on, and why it is listed \emph{before} the analysis questions.
\end{alertbox}

\section{Scales, and What Each Permits}

Stevens' four scale types~\cite{stevens1946theory} determine which operations
are meaningful.

\begin{center}
\small
\begin{longtable}{@{}L{2.0cm}L{3.4cm}L{3.6cm}L{4.4cm}@{}}
\toprule
\textbf{Scale} & \textbf{Structure} & \textbf{Example} & \textbf{Meaningful operations} \\
\midrule
\endfirsthead
\toprule
\textbf{Scale} & \textbf{Structure} & \textbf{Example} & \textbf{Meaningful operations} \\
\midrule
\endhead
\bottomrule
\endfoot
Nominal
  & Categories, no order
  & Programming language; bug type
  & Counts, mode, chi-square \\
Ordinal
  & Ordered, unequal or unknown intervals
  & Severity: low / medium / high; a single Likert item
  & Median, percentiles, rank-based tests \\
Interval
  & Equal intervals, arbitrary zero
  & Calendar date; a summed multi-item scale, approximately
  & Mean, standard deviation, most parametric tests \\
Ratio
  & Equal intervals, true zero
  & Response time; lines of code; defect count
  & All of the above, plus ratios: ``twice as fast'' is meaningful \\
\end{longtable}
\end{center}

\begin{conceptbox}[The Likert controversy, settled practically]
Is it legitimate to take the mean of responses coded 1 to 5? Strictly, no: the
distance between ``disagree'' and ``neutral'' is not known to equal that
between ``neutral'' and ``agree''.

\medskip
In practice the field has reached a workable position:

\begin{tight}
  \item A \textbf{single Likert item} is ordinal. Report the median and the
    distribution; analyse with rank-based methods. Do not report a mean of 3.7.
  \item A \textbf{summed or averaged multi-item scale} with several items
    measuring one construct is conventionally treated as interval. Simulation
    work has repeatedly found parametric tests robust to this, provided the
    scale has enough items and the distribution is not badly skewed.
  \item When in doubt, \textbf{do both} and report that the conclusion does not
    depend on the choice. If it does depend on the choice, that is important
    and you must say so.
\end{tight}

\medskip
The design consequence, stated in \chref{survey}, bears repeating: if you
intend to treat a construct as continuous, write several items for it.
\end{conceptbox}

\section{Validity: Five Questions About One Measure}

\begin{center}
\small
\begin{longtable}{@{}L{3.0cm}L{4.4cm}L{6.8cm}@{}}
\toprule
\textbf{Type} & \textbf{Asks} & \textbf{How you establish it} \\
\midrule
\endfirsthead
\toprule
\textbf{Type} & \textbf{Asks} & \textbf{How you establish it} \\
\midrule
\endhead
\bottomrule
\endfoot
Face
  & Does it look right?
  & Judgement. The weakest form, and not evidence --- but its absence is a
    warning \\
Content
  & Does it cover the whole construct?
  & Expert panel against a definition of the construct's domain. A usability
    measure covering only efficiency has poor content validity \\
Construct
  & Does it measure the intended construct and not a neighbouring one?
  & \emph{Convergent}: correlates with other measures of the same thing.
    \emph{Discriminant}: does \emph{not} correlate too highly with measures of
    different things (\chref{sem}) \\
Criterion
  & Does it predict or agree with an accepted standard?
  & \emph{Concurrent}: agrees with a gold standard now.
    \emph{Predictive}: predicts a future outcome \\
Ecological
  & Does it behave the same outside the study setting?
  & Comparison with field data; often only arguable, not demonstrable \\
\end{longtable}
\end{center}

\begin{paperbox}[Questionable measurement practices]
Jessica Kay Flake and Eiko I. Fried, \emph{Measurement Schmeasurement:
Questionable Measurement Practices and How to Avoid Them}, Advances in Methods
and Practices in Psychological Science 3(4), 2020~\cite{flake2020measurement}.

\medskip
\textbf{Why it matters.} It does for measurement what the replication
literature did for analysis: names the specific practices --- using a measure
with no validity evidence, silently modifying items, selecting among several
measured scales after seeing results --- that quietly invalidate studies.

\medskip
\textbf{What to extract.} Their six questions every author should be able to
answer about every measure in their study. Answer all six for each of yours,
in writing, before data collection.
\end{paperbox}

\section{Reliability, and the Alpha Problem}

Reliability is consistency: would the measure give the same answer again? It is
necessary but not sufficient for validity --- a bathroom scale reading three
kilograms heavy is perfectly reliable and invalid.

\begin{center}
\small
\begin{tabular}{@{}L{3.2cm}L{5.0cm}L{6.1cm}@{}}
\toprule
\textbf{Type} & \textbf{Question} & \textbf{Estimated by} \\
\midrule
Test--retest       & Stable over time?
                   & Correlation between two administrations. Needs a gap long
                     enough to prevent recall, short enough that the construct
                     has not changed \\
Internal consistency & Do the items hang together?
                   & Cronbach's $\alpha$~\cite{cronbach1951coefficient},
                     McDonald's $\omega$ \\
Inter-rater        & Do two observers agree?
                   & Cohen's $\kappa$, Krippendorff's $\alpha$
                     (\chref{qualitative}) \\
Parallel forms     & Do two versions agree?
                   & Correlation between forms \\
\bottomrule
\end{tabular}
\end{center}

\begin{alertbox}[What Cronbach's alpha does not tell you]
$\alpha \ge 0.7$ is quoted as though it certified a scale. It does not, for
three reasons that are well established and widely ignored.

\medskip
\textbf{It is not a measure of unidimensionality.} A scale measuring two
distinct constructs can have a high $\alpha$. Alpha tells you the items
correlate, not that they correlate because of \emph{one} underlying thing.
Establishing dimensionality requires factor analysis (\chref{sem}).

\medskip
\textbf{It increases with the number of items.} Twenty mediocre items will
reach 0.9. Adding items is not improving the scale.

\medskip
\textbf{It assumes tau-equivalence} --- that every item relates equally
strongly to the construct --- which is almost never true. When it is false,
$\alpha$ underestimates or overestimates reliability unpredictably.

\medskip
McDonald's $\omega$ makes weaker assumptions and is available in the R packages
\texttt{psych} and \texttt{MBESS}~\cite{revelle2009coefficients}. Report
$\omega$; report $\alpha$ too if you wish, since reviewers expect it, but do not
let it be your only evidence.
\end{alertbox}

\begin{lstlisting}[language=Rlang, caption={Reliability done properly: dimensionality
first, then omega, then alpha for the reviewers who expect it.}]
library(psych)

# 1. Is it one dimension? Parallel analysis before anything else.
fa.parallel(items, fa = "fa")        # how many factors are indicated?
fa(items, nfactors = 1, rotate = "none")   # do all items load on it?

# 2. Reliability. omega_total is the one to report.
omega(items, nfactors = 1)

# 3. Alpha, with the item-dropped column: an item whose removal
#    raises alpha substantially is measuring something else.
alpha(items)
\end{lstlisting}

\section{Deriving Measures from Goals}

The alternative to measuring what is convenient is to derive measures from what
you want to know. The Goal--Question--Metric approach~\cite{basili1994goal} does
this in three levels.

\begin{center}
\small
\begin{tabular}{@{}L{2.0cm}L{12.3cm}@{}}
\toprule
\textbf{Goal}   & Conceptual. \emph{Analyse} the code review process
                  \emph{for the purpose of} improving it \emph{with respect to}
                  defect detection \emph{from the viewpoint of} the team lead
                  \emph{in the context of} a small product team \\
\textbf{Question} & Operational. How many defects are found in review versus
                  after release? Does review coverage vary by module
                  criticality? \\
\textbf{Metric} & Quantitative. Defects found in review per KLOC; proportion of
                  changed files reviewed; time from submission to first comment \\
\bottomrule
\end{tabular}
\end{center}

The value of the goal template is that it forces the viewpoint and context to
be stated. The same metric means different things to a team lead and to a
developer, and a metric collected without a viewpoint tends to be used by
whoever finds it convenient --- usually as a performance measure, which is how
measurement programmes get subverted.

\begin{worked}[Operationalising ``developer productivity'']
A candidate for the least valid construct in common use. Watch it fail four
times.

\medskip
\textbf{Attempt 1: lines of code per week.} Rewards verbosity, punishes
deletion and reuse. A developer who removes 500 lines of duplication has
negative productivity. \emph{Construct validity: absent.}

\medskip
\textbf{Attempt 2: commits per week.} Rewards splitting work into small
commits, which is good practice --- and therefore measures commit style, not
output. \emph{Confounded with a practice you may also be studying.}

\medskip
\textbf{Attempt 3: story points completed.} Points are estimated by the team,
often by the same people being measured, and inflate under measurement. A
measure that changes when you start using it is not measuring a stable
property. \emph{Reactivity.}

\medskip
\textbf{Attempt 4: features delivered per quarter.} Better --- closer to value
--- but features differ enormously in size, and this ignores maintenance,
review, mentoring and incident response, which is most of what senior
developers do. \emph{Content validity: poor, systematically biased against
senior staff.}

\medskip
\textbf{The honest conclusion.} ``Productivity'' as a single number for an
individual developer may not be a valid construct at all. The defensible moves
are: measure something narrower and well-defined (cycle time for a change of a
stated class), measure at the team level where the output is more meaningful,
or reframe the research question so it does not require the construct.

\medskip
Discovering this in week 10 is a good outcome. Discovering it at the viva is
not.
\end{worked}

\begin{pitfall}[Proxy measures: three questions before you accept one]
Most computing measures are proxies. That is acceptable if you can answer:

\begin{tightnum}
  \item \textbf{What is the evidence the proxy tracks the construct?} A
    citation, or a validation in your own data. Not plausibility.
  \item \textbf{What else does it track?} Cyclomatic complexity correlates with
    lines of code, so ``complexity predicts defects'' may be ``size predicts
    defects''.
  \item \textbf{Does it change when people know it is measured?} If yes,
    Goodhart's law applies: it will stop measuring the construct precisely when
    it starts mattering.
\end{tightnum}
\end{pitfall}

\begin{traditions}{measurement}
\tradEMP{Validity and reliability are technical properties with tests. Report
both for every measure, and treat unvalidated measures as a threat.}
\tradDSR{Evaluation criteria are measures too, and inherit every problem here.
An artefact evaluated on an invalid measure has not been evaluated.}
\tradQUAL{Rejects the framing: the analogue is trustworthiness --- credibility,
transferability, dependability, confirmability --- established by audit trail
and member checking rather than by coefficients (\chref{qda}).}
\tradFORM{Measurement is definitional: the quantity is what the definition
says. The validity question reappears as whether the formal definition captures
the informal notion.}
\end{traditions}

\begin{lab}[Audit your own measures]
For every measure in your alignment table (\chref{questions}):

\begin{tightnum}
  \item Name the construct and the measure separately.
  \item State the scale type, and check that your planned analysis permits it.
  \item Answer Flake and Fried's question: what is the evidence this measures
    the construct? If the answer is ``it seems reasonable'', mark it red.
  \item For multi-item scales, state how you will report reliability ---
    dimensionality check, then $\omega$.
  \item For each proxy, answer the three questions above.
  \item Write a GQM goal statement for your main measure, including viewpoint
    and context.
\end{tightnum}

\medskip
Any measure still marked red needs either validation work or replacement. This
is portfolio piece P4.
\end{lab}

\begin{checkpoint}
\begin{tightnum}
  \item A scale of 20 items has $\alpha = 0.91$. What does this establish, and
    what does it not?
  \item Why can a measure be perfectly reliable and completely invalid? Give a
    computing example.
  \item You report the mean of a single 5-point Likert item. What is the
    objection, and what would you report instead?
  \item Give a construct in your area where measurement reactivity would be a
    serious threat, and say how you would detect it.
\end{tightnum}
\end{checkpoint}

\begin{chaptersummary}
\begin{tight}
  \item Analysis cannot repair measurement, and bad measures produce
    respectable-looking output. \S14.3(B) assesses validity explicitly.
  \item Scale type determines permissible operations. Single Likert items are
    ordinal; multi-item scales are treated as interval, which is why constructs
    need several items.
  \item Five validity questions; construct validity, split into convergent and
    discriminant, is the one that matters most and is checked least.
  \item Reliability is necessary and not sufficient. Cronbach's alpha does not
    establish unidimensionality, rises with item count, and assumes
    tau-equivalence. Check dimensionality, then report McDonald's omega.
  \item GQM derives metrics from goals with a stated viewpoint and context,
    instead of measuring what is available.
  \item Some constructs in common use --- individual developer productivity
    among them --- may not be validly measurable. Narrow the construct, raise
    the level, or change the question.
  \item For every proxy: what is the evidence, what else does it track, and
    does it change under observation?
\end{tight}
\end{chaptersummary}

\begin{reviewq}
  \item Flake and Fried argue that unvalidated measures are as damaging as
    p-hacking, and far less discussed. Assess this claim for computing
    research, and identify which of our subfields is most exposed.
  \item Goodhart's law says a measure that becomes a target ceases to be a good
    measure. What follows for research that studies organisations already using
    the metric you intend to measure?
  \item Construct a case for the position that ``developer productivity'' is a
    coherent construct that we simply measure badly, and then the case that it
    is incoherent. Which do you find stronger, and what evidence would decide?
  \item \reg{PhD Regs 2024}{\S14.3(B)} lists reliability and validity as
    separate criteria. Design a dissertation-level demonstration that a new
    measure in your area satisfies both, and estimate what it would cost in
    time.
\end{reviewq}
